%% ch10 --- Stretch and fold: strange attractors
%%
%% Written September 2026. Every number the reader cannot do in their head
%% comes from code/measure/ch10.py, which imports the SAME rules the listings
%% print (code/ch10/). The tent map and Henon's map use only +, - and *, so
%% every trajectory behind a number on this page is bit-for-bit the same on
%% every machine; the one logarithm in a loop (the Lyapunov exponent) is only
%% summed and is printed to two decimals.
%%
%% WHAT THIS CHAPTER MAY NOT SAY. It shows that the Henon attractor has zero
%% area and layers within layers, and says only that its dimension lies
%% between one and two: MEASURING that dimension (box counting) is Chapter
%% 11's payoff. It uses the Lyapunov exponent as Chapter 7 defines it and the
%% attractor idea Chapters 4 and 9 plant; it does not re-derive the horizon
%% formula (Chapter 7) or the statistics of an attractor (Chapter 9). The
%% floating-point collapse of the tent map after about fifty kneads is
%% Chapter 16's, which is why no run here is longer than fourteen kneads.
%%
%% Twin: chapters/pl/ch10-strange-attractors.tex. Section for section, macro
%% for macro, number for number; tools/parity.py compares them.

\chapter{Stretch and fold: strange attractors}
\label{ch:strange}
\index{strange attractor}\index{stretch and fold}

\noindent
In Chapter~\ref{ch:lorenz} two runs of Lorenz's equations that began a hair
apart drifted until they had nothing to do with each other, and yet neither
ever left the butterfly-shaped region it started in. Put those two facts side
by side and they look impossible. If nearby states are pushed apart at every
step, why do they not end up at opposite ends of the universe? And if
everything stays inside a small region, where does all that separation go?

The answer is something every baker does with their hands. This chapter
shows it on a strip of dough, then on a two-line rule the astronomer Michel
Hénon wrote down in 1976, and then follows it to a shape with no area, no end
to its length, and layers inside its layers however closely you look.

\begin{outcomes}
\outcome{explain how motion can stay inside a region for ever and still be
  sensitive to its start, using the words stretch and fold}
\outcome{iterate the tent map and Hénon's map in Python and read what they
  print}
\outcome{work out how much one step of Hénon's map multiplies an area, and
  check it by measuring}
\outcome{say what an attractor, its basin and a strange attractor are}
\outcome{describe what zooming in shows on a strange attractor, and how that
  differs from zooming in on a smooth curve}
\outcome{describe how a Poincaré section turns a flow into a map}
\end{outcomes}

\section{Stretch, then fold}
\label{sec:ch10-knead}
\index{kneading}\index{tent map}

\noindent
Put a lump of dough on a board and press two raisins into it side by side.
Now knead: roll the dough out to twice its length, then fold the far half
back on top of the near half. Do it again, and again. After a minute the
raisins are nowhere near each other, and neither has left the board.

A knead is two moves, and they are the whole of this chapter. The
\emph{stretch} pulls every pair of neighbours further apart. The \emph{fold}
puts the stretched dough back into the space it came from. Stretching alone
would carry the raisins off the board; folding alone would do nothing to
them. Together they do what neither can.

To turn that into arithmetic, flatten the dough into a strip running from
$0$ to $1$ and let $x$ be a raisin's position along it. Stretching sends $x$
to $2x$. Anything now past $1$ is folded back, so a point that stuck out to
$2x$ lands at $2 - 2x$, as far short of the end as it had gone past it:
\[ x_{n+1} = 2x_n \ \text{if } 2x_n \le 1, \qquad
   x_{n+1} = 2 - 2x_n \ \text{otherwise}, \]
where $x_n$ is the raisin's position after $n$ kneads. Its graph climbs in a
straight line from $0$ to $1$ and comes straight back down, like a tent seen
from the side, so it is called the \emph{tent map}.

\begin{think}
Two raisins start at $\num{0.2345}$ and $\num{0.2355}$, a thousandth of the
strip apart. While both stay on the same side of the fold, every knead
doubles the gap between them. If it went on doubling, after how many kneads
would the gap be longer than the whole strip? What must happen instead?
\end{think}
\reveal{After ten: $\num{0.001} \times 2^{10} = \num{1.024}$, longer than the
strip. It cannot happen: two raisins on a strip of length one are never more
than one apart, so before then the fold must catch one raisin and not the
other.}

\pyregion{code/ch10/kneading.py}{knead}{Kneading a strip of dough: stretch, then fold}{lst:ch10-knead}

\transcript{ch10-kneading}

\noindent
The gap doubles nine times in a row, to \num{0.512}. At the tenth knead one
raisin is folded and the other is not, and the gap is \val{ch10.knead.ten}.
One knead later they are only \val{ch10.knead.eleven} apart, by chance:
scattered raisins can be any distance apart, and from now on knowing where
one is tells you nothing about the other. Every knead multiplies every short gap by
exactly $2$, so this rule's Lyapunov exponent, in the sense of
Chapter~\ref{ch:lyapunov}, is $\ln 2 \approx \val{ch10.knead.lambda}$ per
knead.

\begin{trapbox}
\textbf{\enquote{If a system can never leave a small region, it cannot be very
sensitive to its start.}} The region limits how \emph{big} an error can get,
not how fast it grows. A tiny error doubles and doubles until it is as large
as the region, and from then on the two runs are as unrelated as two raisins
dropped on the board at random. Staying in the bowl and being predictable are
different promises, and the fold is what lets a system keep the first
without the second.
\end{trapbox}

\begin{lab}{l10_01_knead}{A strip of dough}
Write \code{knead}, one stretch and one fold. Then write
\code{kneads\_until\_apart}, which counts the kneads before two raisins that
start a small gap apart are more than a tolerance apart. For raisins a
thousandth apart and a tolerance of \num{0.1} the answer is seven. Start them
ten times closer and see how few extra kneads that buys.
\end{lab}

\section{Hénon's two-line map}
\label{sec:ch10-henon}
\index{Hénon map}\index{Hénon, Michel}

\noindent
The strip is one number. The Lorenz system of Chapter~\ref{ch:lorenz} has
three, and it stretches and folds as a trajectory swings round a lobe, is
pulled away from its neighbours and is thrown back into the middle. In 1976
Michel Hénon, an astronomer at the observatory in Nice, published a rule,
about as simple as a rule can be, that does the same thing to a plane
\dash{} so simple that a computer can run it for millions of steps. It takes
a point $(x, y)$, a state of two numbers, to a new point:
\[ x_{n+1} = 1 - a\,x_n^{2} + y_n, \qquad y_{n+1} = b\,x_n, \]
with $a = \val{ch10.a}$ and $b = \val{ch10.b}$. In words: the new $x$ is one,
minus $a$ times the square of the old $x$, plus the old $y$; the new $y$ is
the old $x$ shrunk to three tenths of its size.

Hénon built it from three moves. A \emph{bend} moves every point up or down
by $1 - a x^{2}$, a hump that sends $x$ and $-x$ to the same height, so a
straight line across the plane is creased into an arch: that is the fold. A
\emph{squeeze} presses the plane sideways to three tenths of its width. A
\emph{swap} exchanges the two numbers, so the next bend works across the
last one.

\begin{think}
Start at the origin, $(0, 0)$. Where is the point after one step of Hénon's
map, and after two?
\end{think}
\reveal{After one step, $(1, 0)$. After two, $x = 1 - \num{1.4} \times 1 +
0 = -\num{0.4}$ and $y = \num{0.3} \times 1$, so the point is
$(-\num{0.4}, \num{0.3})$.}

\pyregion{code/ch10/henon.py}{rule}{Hénon's map: two numbers, one line of arithmetic each}{lst:ch10-henon}

\transcript{ch10-henon}

\noindent
The listing carries on where your hand stopped. The point jumps about without
settling, repeating or running away: however long you let it go, $x$ stays
between \val{ch10.attractor.left} and \val{ch10.attractor.right}. The rule
uses only multiplication, addition and subtraction, so every digit printed is
the same on every computer.

\section{One shape for every start}
\label{sec:ch10-cloud}
\index{attractor}\index{basin of attraction}

\noindent
One start is one story. Take many: a grid of forty by forty,
\val{ch10.cloud.n} starting points spread evenly over the square in which
both $x$ and $y$ run from $-\num{0.5}$ to $\num{0.5}$.

\begin{think}
Run every one of these points through Hénon's map for thirty steps. Where do
you expect the cloud to be: still spread over its square, all gathered at a
single point, or somewhere else?
\end{think}
\reveal{Somewhere else. A few points leave for good, and all the rest lie on
one thin, bent set \dash{} the set a single start draws if you run it long
enough. They do not gather at a point, because the stretching keeps them
apart along that set.}

\pyregion{code/ch10/cloud.py}{cloud}{A cloud of starts, and whether each lands on the attractor}{lst:ch10-cloud}

\transcript{ch10-cloud}

\noindent
To decide whether a point is on that set, the program first runs one start
for \val{ch10.cloud.reference} steps and marks every little square of side
\val{ch10.cloud.cell} that it visits. At the start only \val{ch10.cloud.zero}
of the \val{ch10.cloud.n} points sit in a marked square. From step
\val{ch10.cloud.settled} on, all \val{ch10.cloud.kept} that remain do, and
\val{ch10.cloud.left} have left for good, running off past ten and on
towards infinity.

That is what the word \emph{attractor} means: a set of states that the
motion settles onto and never leaves, and that pulls in the states around it.
The resting point of Chapter~\ref{ch:motion}'s damped pendulum is the
simplest attractor; the butterfly of Chapter~\ref{ch:lorenz} is another.
Every start that ends up on an attractor belongs to its \emph{basin}, as
every raindrop that reaches a lake fell in the lake's basin. Hénon's attractor has a basin with an edge, and starts outside
it run off to infinity.

\plotfig{ch10-attractor}{Hénon's attractor (blue) on its basin. Every start
in the white region ends up on the attractor; every start in the grey region
leaves for good. The top edge of the listing's square reaches into the grey,
and that is where the starts that left began.}{the attractor and its basin}

\section{Squeezed, not stopped}
\label{sec:ch10-area}
\index{area}\index{determinant}\index{dissipation}

\noindent
The cloud began as a square with an area and ended on a set as thin as a
line. Hénon's rule says exactly how fast the area goes. Take a tiny square of
side $h$ with a corner at $(x, y)$. One step moves its two sides, the arrows
$(h, 0)$ and $(0, h)$, to two new arrows. Nudging $x$ by $h$ changes the new
$x$ by about $-2ax\,h$, because $(x + h)^{2} = x^{2} + 2xh + h^{2}$ and
$h^{2}$ is tiny next to $h$, and it changes the new $y$ by $bh$. Nudging $y$
by $h$ changes the new $x$ by $h$ and the new $y$ not at all. So the square
becomes a slanted parallelogram with sides
\[ (p, q) = (-2ax\,h,\ b\,h), \qquad (r, s) = (h,\ 0). \]
The area of a parallelogram with sides $(p, q)$ and $(r, s)$ is the size of
$ps - qr$, a number called the \emph{determinant}; for a rectangle with sides
$(2, 0)$ and $(0, 3)$ it gives $2 \times 3 - 0 \times 0 = 6$. Here it is
$(-2ax\,h) \times 0 - (b\,h) \times h = -b\,h^{2}$, so the area $h^{2}$
becomes $b\,h^{2}$. Every step multiplies every area by $b = \val{ch10.b}$,
wherever the square was: the position $x$, which decides how much the square
is stretched, is multiplied by zero and drops out.

\begin{think}
A patch of starting points has an area of $1$. What is its area after ten
steps? Does that mean its points slow down and come to rest, as the damped
pendulum of Chapter~\ref{ch:motion} did?
\end{think}
\reveal{About \val{ch10.area.ten}, which is $\num{0.3}^{10}$: six millionths
of what it was. And no, the points do not slow down at all. The area is
squeezed out of the patch; the motion is not.}

\begin{trapbox}
\textbf{\enquote{If the area keeps shrinking, the motion must be dying out.}}
A damped pendulum loses area because everything is pulled towards one resting
point. Hénon's map squeezes hard in one direction while it stretches in
another, so a patch becomes a longer, thinner ribbon whose points go on moving
for ever. Shrinking area says that the motion ends up on something thin, not
that it stops.
\end{trapbox}

\noindent
Now measure it. A square a tenth on a side becomes sixteen thousand points
round its edge; the map moves each one, and the surveyor's shoelace formula
measures the area inside the new outline.

\pyregion{code/ch10/area.py}{area}{Measuring what the map does to an area}{lst:ch10-area}

\transcript{ch10-area}

\noindent
The area falls by the factor the determinant promised, to five decimals,
down to \val{ch10.area.six} of the start after six steps. The outline does
the opposite: after first getting shorter, as the squeeze flattens the square,
it keeps growing, and after six steps it is \val{ch10.outline.six} times as
long, wrapped round less than a thousandth of the area.

\plotfig{ch10-square-folds}{A square of starting points, coloured by where
along $x$ each point began, after up to five steps. One step bends it into a
crescent, the next stretches it into a thin strip, and by the third the strip
is folded back on itself and its colours lie in layers.}{a square stretched
and folded}

\noindent
How hard is the stretching? Carry a tiny arrow along one long run, moving it
at each step as the sides of the square were moved. Measure how much longer
it got, add the logarithm of that to a running total, and shrink it back to
length one. The average of the logarithms is the Lyapunov exponent of
Chapter~\ref{ch:lyapunov}, now for a map of the plane.

\pyregion{code/ch10/stretch.py}{arrow}{The average stretch of a tiny arrow}{lst:ch10-arrow}

\transcript{ch10-stretch}

\noindent
The first lines come from the other half of the file: two starts $10^{-10}$
apart, which Python writes as \code{1.0e-10}. Their gap grows until by step
$70$ it is about $2$, and after that it never exceeds \val{ch10.gap.widest},
roughly the width of the attractor. Bounded, and sensitive, like the dough.
The exponent is \val{ch10.lambda} per step: along the attractor a tiny arrow
grows on average by a factor of \val{ch10.stretch} a step, while across it the
map squeezes by \val{ch10.squeeze}. Multiply the two and you get back about
\val{ch10.b}, the area factor. Stretch one way, squeeze the other, fold to
fit: that is the whole recipe.

\begin{lab}{l10_02_henon}{One step, and its area}
Write \code{step}, one step of Hénon's map, and \code{area\_factor}, which
moves three corners of a tiny square one step, reads off the two new sides as
arrows $(p, q)$ and $(r, s)$, and returns the size of $ps - qr$ divided by the
old area. Check that it gives \num{0.3} wherever you put the square, and that
changing $b$ changes it while changing $a$ does not.
\end{lab}

\section{Layers within layers}
\label{sec:ch10-zoom}
\index{strange attractor!layers}\index{fractal}

\noindent
Drawn at the size of a page, Hénon's attractor looks like a few curved lines.
Look closer.

\begin{think}
Draw any smooth curve with a pen and zoom in on a small piece of it, again
and again. What does the piece look like in the end? Now guess what you will
see if you zoom in on one of the attractor's lines.
\end{think}
\reveal{A smooth curve looks more and more like a straight line. The
attractor's line does not: under the magnifier it splits into several
parallel lines, and each of those, magnified again, splits again.}

\plotfig{ch10-zooms}{Zooming into Hénon's attractor. Each red box is shown
enlarged in the next panel. What looked like one line is a bundle of lines,
and each line of the bundle is itself a bundle.}{the attractor, zoomed twice}

\noindent
The layers come straight from the recipe. Every step stretches the set,
folds it over and squeezes it to fit, so each fold lays a thinner copy of
everything beside what was there, and the copy carries its own layers inside
it. Two consequences sound absurd until you have watched them happen.
The attractor has no area, since every area is multiplied by \val{ch10.b} at
every step, for ever. And its length has no end, since it is traced by a curve
stretched by about \val{ch10.stretch} at every step, for ever.

\begin{trapbox}
\textbf{\enquote{Zoom in far enough on any curve and it turns into a straight
line.}} That is true of the smooth curves you draw with a pen and of the
graphs of the formulas you met at school. It is not true here. However far you
zoom into Hénon's attractor, each line splits into a bundle: the structure has
no smallest scale.
\end{trapbox}

\noindent
An attractor with this layered structure, on which the motion is chaotic, is
called a \emph{strange attractor}, a name given by David Ruelle and Floris
Takens in 1971. Lorenz's butterfly is one, and so is Hénon's. Such a set is
more than a line, which would have a finite length, and less than an area,
since it has none: its dimension lies between one and two, and
Chapter~\ref{ch:fractals} measures where.

\begin{worldbox}
Chemical engineers mix thick fluids \dash{} resins, sealants, food pastes
\dash{} with a \emph{static mixer}, a pipe fitted with a row of short twisted
blades and no moving parts. Each blade splits the stream in two, turns the
halves and stacks them: a stretch and a fold, done by the shape of the pipe.
Every blade doubles the number of layers, so twenty blades make roughly a
million, because ten doublings make about a thousand.
\end{worldbox}

\begin{rigourbox}
Stephen Smale turned the baker's knead into mathematics about a decade
before Hénon's map. His
\emph{horseshoe} stretches a square into a long strip, bends it into a
horseshoe and lays it back across the square. The points that stay in the
square for ever form a set with gaps inside gaps, and on it Smale proved that
the motion is as unpredictable as a sequence of coin tosses. What this chapter
shows about Hénon's map at $a = \val{ch10.a}$ and $b = \val{ch10.b}$ is
measured, not proved: nobody has proved that the set the computer draws is not,
say, an extremely long cycle. Michael Benedicks and Lennart Carleson proved in
1991 that maps of Hénon's form do have strange attractors for many other
choices of the two numbers.
\end{rigourbox}

\section{Cutting a flow into a map}
\label{sec:ch10-section}
\index{Poincaré section}

\noindent
Why should a rule with two numbers tell you anything about the Lorenz system,
a flow with three? Picture the butterfly with a pane of glass pushed through
one of its wings. The trajectory pierces the pane again and again, leaving a
dot each time. Forget the path between the crossings and keep only the dots.
Given one dot, the flow decides where the next will be, so the dots obey a
rule of their own: a map from the pane to itself, with two numbers for a
state. This is a \emph{Poincaré section}, after Henri Poincaré, who used it in
his work on the three bodies of Chapter~\ref{ch:mechanics}.

\mermaidfig{ch10-section}{A Poincaré section. The flow's path crosses a pane
again and again; the rule that takes one crossing to the next is a map of
the pane.}{flow, section, map}

\noindent
Hénon's map is not the Lorenz flow's own section, but it was built to behave
like one: between two crossings a flow like Lorenz's stretches a patch of the
pane, squeezes it and folds it, which is what the bend, the squeeze and the
swap do.

\begin{think}
Chapter~\ref{ch:motion} said that a smooth flow needs at least three numbers
in its state to be chaotic. Hénon's map is chaotic with two, and the logistic
map of Chapter~\ref{ch:logistic} with one. Is that a contradiction?
\end{think}
\reveal{No. That rule is about flows, whose paths move continuously and in a
plane can never cross one another. A map jumps from one state to the next, so
nothing stops it folding. Cutting a three-number flow with a pane leaves a
two-number map, and the chaos survives the cut.}

\begin{summarybox}
\begin{itemize}
\item Chaos in a bounded system is two moves repeated: a \textbf{stretch},
  which pulls neighbours apart, and a \textbf{fold}, which puts everything back.
  The \textbf{tent map} kneads a strip and doubles every small gap.
\item \textbf{Hénon's map}, $x_{n+1} = 1 - a x_n^{2} + y_n$ and
  $y_{n+1} = b x_n$, bends, squeezes and swaps a plane.
\item A cloud of starts ends up on one set, the \textbf{attractor}; the starts
  that end there form its \textbf{basin}.
\item Every step multiplies every area by $b$, the \textbf{determinant} of the
  step, yet the motion never stops: the map \textbf{squeezes} one way while it
  \textbf{stretches} the other.
\item Zoomed in, each line of the attractor splits into a bundle, again and
  again. An attractor like that, with chaotic motion on it, is a
  \textbf{strange attractor}: no area, no end to its length.
\item A \textbf{Poincaré section} turns a flow into a map, which is why a
  two-number map can be chaotic although a two-number flow cannot.
\end{itemize}
\end{summarybox}

\canyou

\begin{checkpoint}
\item A raisin sits at $\num{0.3}$ on the strip. Where is it after one knead,
  and after two?
  \answerto{At \num{0.6}, then at \num{0.8}: doubling \num{0.6} gives
  \num{1.2}, past the end, so it folds back to $2 - \num{1.2} = \num{0.8}$.}
\item Two raisins start \num{0.01} apart and stay on the same side of the fold.
  How far apart are they after three kneads, and why can the doubling not go
  on for ever?
  \answerto{\num{0.08}. The strip is only one long, so sooner or later the
  fold catches one raisin and not the other.}
\item Take one step of Hénon's map from the point $(\num{0.5}, 0)$.
  \answerto{$1 - \num{1.4} \times \num{0.25} + 0 = \num{0.65}$ and
  $\num{0.3} \times \num{0.5} = \num{0.15}$: the point $(\num{0.65},
  \num{0.15})$.}
\item A patch of starting points has an area of $2$. What is its area after
  three steps of Hénon's map, and have its points stopped moving?
  \answerto{$2 \times \num{0.027} = \num{0.054}$. No: the patch is squeezed one
  way and stretched another, and its points keep moving.}
\item What is the difference between an attractor and its basin?
  \answerto{The attractor is the set the motion settles onto and never leaves;
  its basin is every start that ends up on it.}
\item What do you see when you zoom in on a smooth curve, and when you zoom in
  on Hénon's attractor?
  \answerto{A smooth curve looks more and more like a straight line. Each line
  of the attractor splits into a bundle of lines, again and again.}
\item What is a Poincaré section, and why does it let a map with a two-number
  state be chaotic?
  \answerto{It keeps only the points where a flow's path crosses a surface;
  the rule from one crossing to the next is a map. The three-number rule is
  about flows, and a map jumps, so it can fold.}
\end{checkpoint}
